% Part: propositional-logic
% Chapter: syntax-and-semantics
% Section: valuations-sat

\documentclass[../../../include/open-logic-section]{subfiles}

\begin{document}

\olfileid[tr]{pl}{syn}{val}

\olsection{\usetoken{P}{valuation} ve Sağlama}

\begin{defn}[!!^{valuation}s]
$\{\True, \False\}$, ``doğru'' ve ``yanlış'' doğruluk değerlerinden oluşan
iki elemanlı küme olsun. $\Lang{L_0}$ için bir \emph{!!{valuation}}, dildeki
!!{propositional variable}s öğelerine $\True$ ya da $\False$ değerlerinden birini
veren, yani $\pAssign{v} \colon \PVar \to \{\True, \False \}$ olan bir
fonksiyon~$\pAssign{v}$'dir.
\end{defn}

\begin{defn}
  Bir !!a{valuation}~$\pAssign{v}$ verildiğinde, değerlendirme fonksiyonunu
  $\pValue{v} \colon \Frm[L_0] \to \{\True, \False \}$ tümevarımlı olarak
  şöyle tanımlayın:
  \begin{align*}
    \iftag{prvFalse}{\pValue{v}(\lfalse) & = \False; \\}{}
    \iftag{prvTrue}{\pValue{v}(\ltrue) & = \True; \\}{}
    \pValue{v}(\Obj p_n) & = \pAssign{v}(\Obj p_n); \\
    \iftag{prvNot}{\pValue{v}(\lnot !A) & = \begin{cases}
        \True & \text{eğer } \pValue{v}(!A) = \False;\\
        \False & \text{aksi hâlde.}
      \end{cases} \\ }{}
    \iftag{prvAnd}{\pValue{v}(!A \land !B) & = \begin{cases}
        \True &
        \text{eğer $\pValue{v}(!A) = \True$ ve $\pValue{v}(!B) = \True$;}\\
        \False &
        \text{eğer $\pValue{v}(!A) = \False$ veya $\pValue{v}(!B) = \False$.}
      \end{cases}\\}{}
    \iftag{prvOr}{\pValue{v}(!A \lor !B) & = \begin{cases}
        \True &
        \text{eğer $\pValue{v}(!A) = \True$ veya $\pValue{v}(!B) = \True$;}\\
        \False &
        \text{eğer $\pValue{v}(!A) = \False$ ve $\pValue{v}(!B) = \False$.}
      \end{cases}\\}{}
    \iftag{prvIf}{\pValue{v}(!A \lif !B) & = \begin{cases}
        \True &
        \text{eğer $\pValue{v}(!A) = \False$ veya $\pValue{v}(!B) = \True$;}\\
        \False &
        \text{eğer $\pValue{v}(!A) = \True$ ve $\pValue{v}(!B) = \False$.}
      \end{cases}\\}{}
    \iftag{prvIff}{\pValue{v}(!A \liff !B) & = \begin{cases}
        \True &
        \text{eğer $\pValue{v}(!A) = \pValue{v}(!B)$;}\\
        \False &
        \text{eğer $\pValue{v}(!A) \neq \pValue{v}(!B)$.}
      \end{cases}}{}
\end{align*}
\end{defn}

\begin{explain}
Bu maddeler aşağıdaki doğruluk tablolarına karşılık gelir:
\begin{center}
\iftag{prvNot}{
  \begin{tabular}{|c||c|} \hline
    $!A$ & $ \lnot !A$ \\
    \hline \hline
    $\True$ & $\False$ \\
    $\False$ & $\True$ \\
    \hline
  \end{tabular}
}{}
\iftag{prvAnd}{
  \begin{tabular}{|cc||c|} \hline
    $!A$ & $!B$ & $!A \land !B$ \\
    \hline \hline
    $\True$ & $\True$ & $\True$ \\
    $\True$ & $\False$ & $\False$ \\
    $\False$ & $\True$ & $\False$ \\
    $\False$ & $\False$ & $\False$ \\
    \hline
  \end{tabular}
}{}
\iftag{prvOr}{
  \begin{tabular}{|cc||c|} \hline
    $!A$ & $!B$ & $!A \lor !B$ \\
    \hline \hline
    $\True$ & $\True$ & $\True$ \\
    $\True$ & $\False$ & $\True$ \\
    $\False$ & $\True$ & $\True$ \\
    $\False$ & $\False$ & $\False$ \\
    \hline
  \end{tabular}
}{}%
\iftag{notprvNot,notprvAnd,notprvOr,notprvIf}{}{\\[1em]}
\iftag{prvIf}{
  \begin{tabular}{|cc||c|} \hline
    $!A$ & $!B$ & $!A \lif !B$ \\
    \hline \hline
    $\True$ & $\True$ & $\True$ \\
    $\True$ & $\False$ & $\False$ \\
    $\False$ & $\True$ & $\True$ \\
    $\False$ & $\False$ & $\True$ \\
    \hline
  \end{tabular}
}{}
\iftag{prvIff}{
  \begin{tabular}{|cc||c|} \hline
    $!A$ & $!B$ & $!A \liff !B$ \\
    \hline \hline
    $\True$ & $\True$ & $\True$ \\
    $\True$ & $\False$ & $\False$ \\
    $\False$ & $\True$ & $\False$ \\
    $\False$ & $\False$ & $\True$ \\
    \hline
  \end{tabular}
}{}
\end{center}
\end{explain}

\begin{prob}
$\Lang{L_0}$'a, değerlendirmesi
\begin{gather*}
  \pValue{v}(\diamondsuit( !A , !B , !C ) ) = \begin{cases}
    \pValue{v}( !B ) &
    \text{eğer $\pValue{v}(!A) = \True$;}\\
    \pValue{v}( !C ) &
    \text{eğer $\pValue{v}(!A) = \False $}
  \end{cases}
\end{gather*}
ile verilen üçlü bir $\diamondsuit$ bağlacı eklediğimizi düşünün. Bu bağlacın
doğruluk tablosunu yazın.
\end{prob}

\begin{thm}[Yerel Belirlenim]
\ollabel{thm:LocalDetermination} $\pAssign{v_1}$ ve $\pAssign{v_2}$,
!!{valuation}s olsun ve $!A$'da geçen !!{propositional
variable}s üzerinde uzlaşsın; yani $\Obj p_n$, $!A$'da her geçtiğinde
$\pAssign{v_1}(\Obj p_n) = \pAssign{v_2}(\Obj p_n)$ olsun. O zaman
$\pValue{v_1}$ ile $\pValue{v_2}$ de $!A$ üzerinde uzlaşır; yani
$\pValue{v_1}(!A) = \pValue{v_2}(!A)$ olur.
\end{thm}

\begin{proof}
$!A$ üzerinde tümevarımla.
\end{proof}

\begin{defn}[Sağlama]
\ollabel{defn:satisfaction} !!a{formula}~$!A$ ve
  !!a{valuation}~$\pAssign{v}$ için \emph{$!A$'nın $\pAssign{v}$ tarafından
  sağlanması} kavramını, $\pSat{v}{!A}$, tümevarımlı olarak aşağıdaki gibi
  tanımlayabiliriz.
  (``$\pSat{v}{!A}$ değil'' demek için $\pSat/{v}{!A}$ yazarız.)
\begin{enumerate}
\tagitem{prvFalse}{%
  \indcase{!A}{\lfalse}{$\pSat/{v}{\indfrm}$.}}{}

\tagitem{prvTrue}{%
  \indcase{!A}{\ltrue}{$\pSat{v}{\indfrm}$.}}{}

\item \indcase{!A}{\Obj p_i}{$\pSat{v}{\indfrm}$ ancak ve ancak
  $\pAssign{v}(\Obj p_i) = \True$ ise.}

\tagitem{prvNot}{%
  \indcase{!A}{\lnot !B}{$\pSat{v}{\indfrm}$ ancak ve ancak
    $\pSat/{v}{!B}$ ise.}}{}

\tagitem{prvAnd}{%
  \indcase{!A}{(!B \land !C)}{$\pSat{v}{\indfrm}$ ancak ve ancak
    $\pSat{v}{!B}$ ve $\pSat{v}{!C}$ ise.}}{}

\tagitem{prvOr}{%
  \indcase{!A}{(!B \lor !C)}{$\pSat{v}{\indfrm}$ ancak ve ancak
    $\pSat{v}{!B}$ veya $\pSat{v}{!C}$ (ya da her ikisi) ise.}}{}

\tagitem{prvIf}{%
  \indcase{!A}{(!B \lif !C)}{$\pSat{v}{\indfrm}$ ancak ve ancak
    $\pSat/{v}{!B}$ veya $\pSat{v}{!C}$ (ya da her ikisi) ise.}}{}

\tagitem{prvIff}{%
  \indcase{!A}{(!B \liff !C)}{$\pSat{v}{\indfrm}$ ancak ve ancak
    $\pSat{v}{!B}$ ile $\pSat{v}{!C}$'nin ikisi de geçerliyse veya
    $\pSat{v}{!B}$ ile $\pSat{v}{!C}$'nin hiçbiri geçerli değilse.}}{}
\end{enumerate}
$\Gamma$ bir !!{formula}s kümesiyse $\pSat{v}{\Gamma}$ ancak ve ancak her
$!A \in \Gamma$ için $\pSat{v}{!A}$ ise geçerlidir.
\end{defn}

\begin{prop}\ollabel{prop:sat-value}
  $\pSat{v}{!A}$ ancak ve ancak $\pValue{v}(!A) = \True$ ise geçerlidir.
\end{prop}

\begin{proof}
  $!A$ üzerinde tümevarımla.
\end{proof}

\begin{prob}
  \olref[pl][syn][val]{prop:sat-value}'ı kanıtlayın.
\end{prob}

\end{document}
